\section{Results}
\label{sec:results}

\textbf{Qwen and DeepSeek evaluations}
\cref{tab:main-results} shows the main Qwen and DeepSeek matrix. Spatial reasoning is essentially unsolved in this setting: the best audited spatial result is 22.0\%, only two percentage points above five-way chance. Affordance results are stronger, but the subcategory analysis in \cref{tab:subcat} shows that this strength is mostly due to object blockage/reachability questions rather than stable grasp reasoning.

\begin{table}[htbp]
  \centering
  \resizebox{\columnwidth}{!}{%
  \begin{tabular}{@{}lcc@{}}
    \toprule
    \textbf{Model} & \textbf{Spatial ZS} & \textbf{Spatial CoT} \\
    \midrule
    Qwen2.5-VL-3B & \textbf{40/200, 20.0\%} (0.18s) & \textbf{39/200, 19.5\%} (4.90s) \\
    DeepSeek-VL2-tiny & \textbf{39/200, 19.5\%} (0.17s) & \textbf{44/200, 22.0\%} (0.74s) \\
    \midrule
    \midrule
    \textbf{Model} & \textbf{Affordance ZS} & \textbf{Affordance CoT} \\
    \midrule
    Qwen2.5-VL-3B & 131/200, 65.5\% (0.46s) & 128/200, 64.0\% (3.40s) \\
    DeepSeek-VL2-tiny & \textbf{144/200, 72.0\%} (0.14s) & \textbf{43/200, 21.5\%} (0.27s) \\
    \bottomrule
  \end{tabular}%
  }
  \caption{Main audited 200-example results. Cells report correct/total, accuracy, and average response time. Bold values mark the near-chance spatial results and the DeepSeek affordance zero-shot/CoT contrast.}
  \label{tab:main-results}
\end{table}

\textbf{Effect of Prompting Strategy: zero-shot vs. CoT.}
CoT changes behavior but does not reliably improve accuracy. For Qwen2.5-VL-3B, CoT slightly hurts both spatial and affordance performance: $-0.5$ percentage points on spatial and $-1.5$ on affordance. DeepSeek-VL2-tiny gains only $+2.5$ points on spatial, but collapses on affordance from 72.0\% zero-shot to 21.5\% CoT. In the paired DeepSeek affordance comparison, 114 examples are correct in zero-shot but wrong under CoT, while only 13 move from wrong to correct. This is the clearest evidence that reasoning traces can introduce systematic errors rather than merely exposing latent reasoning.

\begin{table}[htbp]
  \centering
  \resizebox{\columnwidth}{!}{%
  \begin{tabular}{@{}llccc@{}}
    \toprule
    \textbf{Model} & \textbf{Prompt} & \textbf{Directional orient.} & \textbf{Depth/order} & \textbf{Cross-view} \\
    \midrule
    Qwen2.5-VL-3B & ZS & \textbf{8/76, 10.5\%} & 24/91, 26.4\% & 8/33, 24.2\% \\
    Qwen2.5-VL-3B & CoT & 12/76, 15.8\% & 21/91, 23.1\% & 6/33, 18.2\% \\
    DeepSeek-VL2-tiny & ZS & 14/76, 18.4\% & 17/91, 18.7\% & 8/33, 24.2\% \\
    DeepSeek-VL2-tiny & CoT & 15/76, 19.7\% & 21/91, 23.1\% & 8/33, 24.2\% \\
    \midrule
    \midrule
    \textbf{Model} & \textbf{Prompt} & \textbf{Blockage/reach.} & \textbf{Grasp stability} & \\
    \midrule
    Qwen2.5-VL-3B & ZS & \textbf{105/132, 79.5\%} & 26/68, 38.2\% & \\
    Qwen2.5-VL-3B & CoT & \textbf{107/132, 81.1\%} & 21/68, 30.9\% & \\
    DeepSeek-VL2-tiny & ZS & \textbf{112/132, 84.8\%} & 32/68, 47.1\% & \\
    DeepSeek-VL2-tiny & CoT & \textbf{18/132, 13.6\%} & \textbf{25/68, 36.8\%} & \\
    \bottomrule
  \end{tabular}%
  }
  \caption{Subcategory breakdown for the audited 200-example spatial and affordance slices. Bold values mark the central findings: near-chance directional orientation, strong object blockage/reachability performance, and the DeepSeek CoT collapse on blockage/reachability.}
  \label{tab:subcat}
\end{table}

\begin{table}[t]
  \centering
  \small
  \setlength{\tabcolsep}{4pt}
  \begin{tabular}{@{}lcccc@{}}
    \toprule
    Model & Spatial ZS & Aff. ZS & Spatial CoT & Aff. CoT \\
    \midrule
    SmolVLM-2.2B & \textbf{20\%} & 23\% & 66\% & 47\% \\
    Phi-4-multimodal & \textbf{12\%} & 19\% & 40\% & 44\% \\
    \bottomrule
  \end{tabular}
  \caption{Additional compact-model results.}
  \label{tab:poster}
\end{table}

\textbf{Effect of Depth Information.}
Contrary to expectations, depth information did not show any improvement in model accuracy. Qwen2.5-VL-3B was benchmarked on a dataset of 2084 spatial reasoning examples, which were augmented by concatenating the depth image to the side of the original image. Evaluated with zero-shot prompting, accuracy actually decreased slightly to 18.3\% on RGB-depth images.

\textbf{Sanity probes results.}
\cref{tab:sanity} summarizes the raw image-understanding sanity probes. Both audited models reject blank images reliably, so the image channel is not entirely ignored. The more informative controls are wrong answers and shuffled images. Qwen supports a deliberately wrong answer in 16/20 cases and rejects only 3/20 shuffled-image controls. DeepSeek rejects all shuffled controls, but its correct-answer explanations are usually classified as unclear. These probes show why answer accuracy alone is not sufficient evidence for grounded robotics reasoning.

\begin{table*}[t]
  \centering
  \small
  \setlength{\tabcolsep}{4pt}
  \begin{tabular}{@{}lp{0.34\linewidth}p{0.23\linewidth}p{0.23\linewidth}@{}}
    \toprule
    Probe & Expected behavior & Qwen2.5-VL-3B & DeepSeek-VL2-tiny \\
    \midrule
    Correct answer & Support correct answer on original image & 12 supports, 8 unclear & \textbf{20 unclear} \\
    Wrong answer & Reject deliberately wrong answer & \textbf{16 supports}, 4 unclear & 3 rejects, 17 unclear \\
    Blank image & Reject image as insufficient & 20 rejects & 20 rejects \\
    Shuffled image & Reject mismatched image & \textbf{3 rejects}, 17 unclear & \textbf{20 rejects} \\
    Overlay description & Describe visual overlays without answering & 20 unclear & 20 unclear \\
    \bottomrule
  \end{tabular}
  \caption{Heuristic labels for 20 examples per probe type. Bold values mark behavior most relevant to the grounding claim: unclear correct-answer support, wrong-answer rationalization, and mismatched-image rejection or failure to reject.}
  \label{tab:sanity}
\end{table*}


\textbf{Human baseline.}
The resulting human performance provides an additional reference point for interpreting model results and for assessing benchmark validity. In particular, it helps determine whether incorrect model predictions stem from limitations in model reasoning or from ambiguities inherent in the benchmark itself. \cref{tab:human} indicates that VLMs outperform human participants in overall affordance understanding, with the most pronounced advantage observed in the grasp stability subcategory. In contrast, humans demonstrate stronger performance in spatial reasoning tasks. However, it is noteworthy that performance in the directional orientation subcategory remains close to chance level (five-way guessing), suggesting that both humans and models struggle with this aspect of spatial reasoning using the Robo2VLM dataset.

\begin{table*}[t]
  \centering
  \small
  \setlength{\tabcolsep}{4pt}
  \begin{tabular}{@{}lccc@{}}
    \toprule
    Spatial total & Directional orientation & Depth/distance ordering & Cross-view correspondence \\
    \midrule
    43.3\% & \textbf{26.9\%} & 52.5\% & 42.0\% \\
    \midrule
    Affordance total & Object blockage/reachability & Grasp stability & \\
    \midrule
    45.0\% & 63.3\% & \textbf{22.7\%} & \\
    \bottomrule
  \end{tabular}
  \caption{Human performance using a randomized 50 data-samples subset. Bold values mark the human categories closest to chance and therefore most indicative of benchmark difficulty or ambiguity.}
  \label{tab:human}
\end{table*}
